% !TeX root = ../main.tex
% Fonti: Lezione 5, pp. 1--6; Lezione 6, p. 1 (esercizi iniziali).
\chapter{Spazi vettoriali e sottospazi}
\fonti{Lezioni 5 e 6}

\section{La struttura di spazio vettoriale}
\begin{definizione}
Sia $K$ un campo. Uno \emph{spazio vettoriale su $K$} è un insieme non
vuoto $V$ munito di un'addizione $V\times V\to V$, $(v,w)\mapsto v+w$, e
di una moltiplicazione per scalari $K\times V\to V$,
$(\lambda,v)\mapsto\lambda v$, per cui, per ogni $u,v,w\in V$ e
$\lambda,\mu\in K$, valgono:
\begin{enumerate}[label=(\arabic*)]
\item $(u+v)+w=u+(v+w)$ e $u+v=v+u$;
\item esiste $0_V\in V$ con $v+0_V=0_V+v=v$;
\item per ogni $v$ esiste $-v$ con $v+(-v)=0_V$;
\item $(\lambda\mu)v=\lambda(\mu v)$ e $1_Kv=v$;
\item $(\lambda+\mu)v=\lambda v+\mu v$ e
      $\lambda(v+w)=\lambda v+\lambda w$.
\end{enumerate}
Gli elementi di $V$ si chiamano \emph{vettori}, quelli di $K$ \emph{scalari}.
\end{definizione}

La prima parte degli assiomi dice che $(V,+)$ è un gruppo abeliano. Ne
seguono l'unicità di $0_V$ e dell'opposto: se $e,e'$ sono neutri, allora
$e=e+e'=e'$; se $v+a=v+b=0_V$, aggiungendo $-v$ a entrambi i membri si
ottiene $a=b$. Le leggi distributive permettono poi di ricavare fatti
usati spesso senza essere riscritti. Infatti
$0_Kv=(0_K+0_K)v=0_Kv+0_Kv$, dunque $0_Kv=0_V$; analogamente
$\lambda0_V=0_V$. Inoltre $(-1_K)v=-v$, perché
$v+(-1_K)v=(1_K-1_K)v=0_V$.

\section{Modelli ed esempi}
I segmenti orientati del piano forniscono l'immagine geometrica della
somma e del prodotto per scalare. Spostando i rappresentanti in modo che
il secondo inizi dove termina il primo si ottiene la regola del triangolo.
Ripetendo lo stesso vettore, $v+v=2v$.
\begin{schema}
\begin{tikzpicture}[scale=.8]
  \coordinate (O) at (0,0); \coordinate (A) at (2.5,.35);
  \coordinate (B) at (3.35,1.65);
  \draw[vettore] (O)--(A) node[midway,below] {$v$};
  \draw[vettore] (A)--(B) node[midway,right] {$w$};
  \draw[vettore] (O)--(B) node[midway,above left] {$v+w$};
  \fill (O) circle (1pt) node[below left] {$O$};
\end{tikzpicture}
\end{schema}

Il campo $K$, con le proprie operazioni, è uno spazio vettoriale su se
stesso. Più in generale, se $K$ è un sottocampo di un campo $L$, allora
$L$ è uno spazio vettoriale su $K$: gli scalari sono presi in $K$ e i
vettori in $L$. Sono esempi $\R$ su $\Q$ e $\C$ su $\R$.
\footnote{Per questo esempio basta la relazione di inclusione tra campi;
non occorre una costruzione dei numeri complessi.}

Per $m\geq1$, l'insieme $K^m$ delle $m$-uple è uno spazio vettoriale con
operazioni eseguite componente per componente:
\begin{align*}
 (a_1,\ldots,a_m)+(b_1,\ldots,b_m)
   &=(a_1+b_1,\ldots,a_m+b_m),\\
 \lambda(a_1,\ldots,a_m)&=(\lambda a_1,\ldots,\lambda a_m).
\end{align*}
Le stesse formule valgono se le $m$ componenti vengono disposte in una
colonna. Ogni assioma si verifica componente per componente nel campo $K$.
In $K[x]$, l'insieme dei polinomi in una indeterminata, somma e prodotto
per scalare si ottengono sommando o moltiplicando i coefficienti. Infine,
le funzioni $\R\to\R$ formano uno spazio vettoriale con
$(f+g)(x)=f(x)+g(x)$ e $(\lambda f)(x)=\lambda f(x)$ per ogni $x$.

\section{Riconoscere un sottospazio}
\begin{definizione}
Un sottoinsieme $W\subseteq V$ è un \emph{sottospazio vettoriale} se, con
le operazioni indotte da $V$, è a sua volta uno spazio vettoriale su $K$.
\end{definizione}
È sufficiente controllare che $W$ sia non vuoto, chiuso per somma e per
prodotto per ogni scalare. Infatti, scegliendo $w\in W$, la chiusura per
scalari dà $0_Kw=0_V\in W$ e $(-1_K)w=-w\in W$; gli altri assiomi sono
ereditati da $V$. Equivalentemente, basta verificare che $W$ sia non
vuoto e che $\lambda u+\mu v\in W$ per tutti $u,v\in W$ e
$\lambda,\mu\in K$.

\begin{esempio}
In $\R^2$ l'insieme $\{(a,b):a+b=1\}$ non è un sottospazio:
$(1,0)$ e $(0,1)$ vi appartengono, mentre la loro somma $(1,1)$ no.
Invece $\{(a,b):a+b=0\}$ lo è. Se $a+b=c+d=0$, allora
$(a+c)+(b+d)=0$, e $\lambda a+\lambda b=\lambda(a+b)=0$.
Geometricamente è una retta passante per l'origine.
\end{esempio}
\begin{esempio}
Il primo quadrante $\{(a,b)\in\R^2:a\geq0,\ b\geq0\}$ è chiuso per
somma, ma non per il prodotto per scalari negativi: per esempio
$-1\cdot(1,0)=(-1,0)$ non vi appartiene.
\end{esempio}
\begin{esercizio}
Si consideri $W=\{(a,b)\in\R^2:a^2-2b=0\}$. Benché $0\in W$, non è
un sottospazio. Infatti $(1,\tfrac12)\in W$, ma
$(1,\tfrac12)+(1,\tfrac12)=(2,1)$ non soddisfa l'equazione:
$2^2-2\cdot1=2\ne0$. Anche la moltiplicazione per $2$ fallisce per lo
stesso vettore. Il grafico $b=a^2/2$ è una parabola, non una retta.
\end{esercizio}

\section{Intersezioni e unioni}
\begin{proposizione}
Se $U,W$ sono sottospazi di $V$, allora $U\cap W$ è un sottospazio.
\end{proposizione}
\begin{proof}
Il vettore nullo appartiene a entrambi. Se $u,v\in U\cap W$, essi
appartengono sia a $U$ sia a $W$; per la chiusura di ciascuno,
$u+v\in U\cap W$. Per $\lambda\in K$ vale analogamente
$\lambda u\in U\cap W$.
\end{proof}
L'unione in generale non è un sottospazio. Prendiamo in $\R^2$
$U=\{(a,0):a\in\R\}$ e $W=\{(0,b):b\in\R\}$. La somma di
$(1,0)\in U$ e $(0,1)\in W$ è $(1,1)\notin U\cup W$. Il disegno
mostra perché serva un'operazione diversa dall'unione insiemistica per
ottenere il sottospazio contenente entrambe le rette.
\begin{schema}
\begin{tikzpicture}[scale=.68]
  \draw[asse] (-2.5,0)--(2.8,0) node[right] {$U$};
  \draw[asse] (0,-2)--(0,2.5) node[above] {$W$};
  \draw[vettore] (0,0)--(1.45,0) node[below] {$(1,0)$};
  \draw[vettore] (0,0)--(0,1.4) node[left] {$(0,1)$};
  \draw[vettore] (0,0)--(1.45,1.4) node[right] {$(1,1)$};
  \draw[dashed,black!45] (1.45,0)--(1.45,1.4)--(0,1.4);
\end{tikzpicture}
\end{schema}

\section{Combinazioni lineari e generatori}
\begin{definizione}
Una \emph{combinazione lineare} dei vettori $v_1,\ldots,v_n\in V$ è un
vettore della forma $\lambda_1v_1+\cdots+\lambda_nv_n$, dove gli
scalari $\lambda_i$ appartengono a $K$. Un insieme $S\subseteq V$ è un
\emph{insieme di generatori} di $V$ se ogni vettore di $V$ si scrive
come combinazione lineare finita di elementi di $S$.
\end{definizione}
In $\R^2$, per esempio,
$(a,b)=a(1,0)+b(0,1)$: le due direzioni degli assi generano il piano.
Per $S=\varnothing$ la sola combinazione lineare disponibile è quella
vuota, che per convenzione vale $0_V$.

\begin{esercizio}
In $\R^2$ siano $v_1=(2,-1)$ e $v_2=(1,3)$. Per vedere se generano
il piano, risolviamo
$a v_1+b v_2=(x,y)$, cioè
$2a+b=x$, $-a+3b=y$. Il determinante dei coefficienti è
$2\cdot3-(-1)\cdot1=7\ne0$, dunque per ogni $(x,y)$ esiste una
soluzione unica: $a=(3x-y)/7$ e $b=(x+2y)/7$. In $\R^3$, invece, i
vettori $(1,1,-1)$ e $(2,0,1)$ non possono generare tutto lo spazio:
ogni loro combinazione dipende da due soli parametri e, per esempio,
$(0,1,0)$ non si può ottenere. Infatti dalle prime due coordinate
$a=1$ e $b=-\tfrac12$, mentre la terza darebbe
$-a+b=-\tfrac32\ne0$.
\end{esercizio}
